\section[Fargues--Fontaine 曲线上的 $G$-丛]{Fargues--Fontaine 曲线上的 $G$-丛（12 月 4 日）}

\srcpage{55}
固定任意（连通的）约化群 $G/E$。

\subsection{$G$-丛}

我们回顾 $G$-丛（又称「$G$-主丛」）的一般概念。

\begin{proposition}
设 $X$ 是 $E$ 上的概形。以下三者自然等价：
\begin{enumerate}[label=(\arabic*),itemsep=2pt,topsep=2pt]
\item （「几何 $G$-主丛」）带有 $X$ 上 $G$-作用的概形 $Y\to X$，使得在 $X$ 上局部
（étale/光滑/fppf/fpqc）存在 $G$-等变同构
\[
Y\cong G\times X.
\]
\item （「上同调 $G$-主丛」）$X_\et$ 上的层 $\cF$ 连同 $G$-作用，使得在 $X_\et$ 上
局部存在 étale 层的 $G$-等变同构 $\cF\cong G$。
\item （「Tannaka $G$-主丛」）正合 $\otimes$-函子 $\Rep_E(G)\to\mathrm{VB}(X)$。
\end{enumerate}
\end{proposition}

\begin{example}
若 $G=\mathrm{GL}_n$，则 $\mathrm{GL}_n$-丛等价于 $X$ 上秩 $n$ 的向量丛。例如要得到
描述 (1)，考虑该 $G$-主丛的基的空间。
\end{example}

\begin{example}
若 $G=\mathrm{Sp}_{2n}$，则这些丛的概念等价于 $X$ 上秩 $2n$ 的向量丛 $E$ 连同其上的
完美交错形式。
\end{example}

\begin{proof}
(1)$\to$(2)：取 $Y\to X$ 的截面层。

(2)$\to$(3)：对 $V\in\Rep_E(G)$ 与上同调 $G$-主丛 $\cF$，$V\times^G\cF$ 是 $X_\et$ 上的
$\cO_X$-模，有限秩局部自由。再由 étale 下降即得 $X$ 上的向量丛。

(3)$\to$(1)：考虑 $\cO(G)$ 及其由左、右平移给出的 $G\times G$-作用。它属于
$\Ind(\Rep_E G)$；事实上它是代数对象。于是把正合 $\otimes$-函子
$F\colon\Rep_E(G)\to\mathrm{VB}(X)$ 作用到 $\cO(G)\in\Ind(\Rep_E G)$ 上，得到
$\operatorname{Alg}(\Ind\mathrm{VB}(X))\subset\operatorname{Alg}(\QCoh(X))$ 中带
$G$-作用的对象。再取 $Y=\Spec F(\cO(G))$。
\end{proof}

\begin{remark}
类似的讨论适用于 adic 空间上的 $G$-主丛。注意 adic 空间上拟凝聚层的概念性质不好，
但向量丛的理论是好的，而我们这里实际只用到向量丛。为此使用概形论的曲线：向量丛的
全空间可以代数地构造（用相对 $\Spec$），然后再解析化。

实践中最好用的描述是正合 $\otimes$-函子那一款。
\end{remark}

\begin{corollary}
$X$ 上的 $G$-主丛由 $H^1(X_\et,G)$ 分类。
\end{corollary}

\subsection{$G$-等晶体}

回忆我们定义了 $E$-等晶体的范畴
\[
\Isoc_E=\left\{(V,\varphi):\
\begin{array}{c}
V\text{ 是有限维 }\breveE\text{-向量空间，}\\
\varphi\colon V\xrightarrow{\sim}V\text{ 是 }\varphi_E\text{-线性的}
\end{array}\right\}.
\]

\begin{definition}
$G$-等晶体是正合 $\otimes$-函子 $\Rep_E(G)\to\Isoc_E$。
\end{definition}

\begin{proposition}
任何 $G$-等晶体都具有形式
\[
V\mapsto(V\otimes_E\breveE,b\sigma)\colon\Rep_E(G)\to\Isoc_E,
\]
其中 $b\in G(\breveE)$。这诱导了双射
\[
G(\breveE)/\sigma\text{-共轭}\cong\{G\text{-等晶体}\}/\sim.
\]
\end{proposition}

\begin{remark}
$G$-等晶体应当理解为「$\Spec\breveE/\sigma^\ZZ$ 上的 $G$-主丛」。
\end{remark}

\begin{proof}
\srcpage{56}只需说明 $\Spec\breveE$ 上所有 $G$-主丛都平凡。这来自 Steinberg 的一个定理：
$H^1_\et(\Spec\breveE,G)=\ast$。（这又基于 $\breveE$ 的上同调维数为 $1$。）
\end{proof}

\begin{definition}
\[
B(G):=G(\breveE)/\sigma\text{-共轭}\cong\{G\text{-等晶体}\}/\sim.
\]
元素记作 $b\in B(G)$，常隐含地取定 $G(\breveE)$ 中的一个代表。\yizhu{$B(G)$ 在文献中
亦称 Kottwitz 集（Kottwitz set）；第 10 讲中的 $B(\mathrm{GL}_n)$ 即其特例。}{A（原文
§11.3 的 Kottwitz 分类）}
\end{definition}

\begin{example}
对 $G=\mathrm{GL}_n$，$B(G)$ 是秩 $n$ 等晶体的同构类集合。它可等同于宽度为 $n$ 的
Newton 多边形的集合。
\end{example}

\subsection{Kottwitz 分类}

Kottwitz 对一般的 $G$ 给出了 $B(G)$ 的组合描述。Newton 多边形有一个推广，但一般还需要
多一点信息。粗略地说，答案由 Newton 多边形（满足某个对称性条件）加上有限多的额外数据
给出（对 $\mathrm{GL}_n$ 不需要这些额外数据）。

\subsubsection{Newton 点}

对任何 $(V,\varphi)\in\Isoc_E$，$V$ 由斜率分解自然地 $\QQ$-分次：
$V=\bigoplus_{\lambda\in\QQ}V^\lambda$。于是得到映射
\[
\mathbb{D}\to\mathrm{GL}_{\breveE}(V),
\]
其中 $\mathbb{D}$ 是特征群为 $X^*(\mathbb{D})=\QQ$ 的 pro-环面（从而
$\Rep_E(\mathbb{D})$ 等价于 $\QQ$-分次 $E$-向量空间的范畴）。

若 $F\colon\Rep_E(G)\to\Isoc_E$ 是正合 $\otimes$-函子，则由 Tannaka 形式系统，对所有
$V\in\Rep_E(G)$ 的相容映射 $\mathbb{D}\to\mathrm{GL}_{\breveE}(F(V))$ 等价于同态
$\mathbb{D}\to G_{\breveE}$。因此，对任何等晶体 $b\leftrightarrow F$（其底函子为
$V\mapsto V\otimes_E\breveE$），得到良定义的共轭类
$\nu(b)\colon\mathbb{D}\to G_{\breveE}$。它可以经由一个环面分解。取 $X=X_*(T)$，其中
$T\subset B\subset G_{\breveE}$。尽管从表述上看不出来，$X_*(T)$ 典范地不依赖于 $T$ 的
选取，从而继承 $\Gamma:=\Gal(\bar E/E)$ 的作用。若 $X^+\subset X$ 是支配余特征的子集，
则存在唯一的代表 $\nu(b)\in(X^+_{\QQ})^\Gamma$。这就是 $b$ 的 Newton 点，而映射
$B(G)\to(X^+_{\QQ})^\Gamma$ 称为 Newton 映射。

\begin{example}
对 $G=\mathrm{GL}_n$ 及通常的 $B,T$，有 $X=X_*(T)=\ZZ^n$，
$X^+=\{(m_1\ge m_2\ge\cdots\ge m_n)\}$。由于 $G$ 分裂，$\Gamma$ 的作用平凡。有
\[
X^+_{\QQ}=\{(\lambda_1\ge\cdots\ge\lambda_n):\lambda_i\in\QQ\}.
\]
对秩 $n$ 的等晶体，这只是记录各斜率。
\end{example}

\begin{example}
对 $\mathrm{GL}_n$，Newton 映射
\[
\nu\colon B(G)\to(X^+_{\QQ})^\Gamma
\]
是单射，但对一般的 $G$ 不然。
\end{example}

\begin{example}
设 $G=T$ 是环面。这等价于数据 $X=X_*(T_E)$ 连同其 $\Gamma$-作用。所以
$B(T)=T(\breveE)/\sigma\text{-共轭}\cong T(\breveE)/(\sigma-1)$。
\end{example}

\begin{proposition}
存在同构 $B(T)\cong X_*(T)_\Gamma$，它对 $E$ 上环面之间的映射函子性。在该同构下，
$\nu\colon B(T)\to(X^+_{\QQ})^\Gamma=X^\Gamma_{\QQ}$ 由对 $\Gamma$ 取平均给出：
\[
\text{「}\ \frac{1}{|\Gamma|}\sum_{\gamma\in\Gamma}\gamma\colon
X_*(T)_\Gamma\to X_*(T)^\Gamma_{\QQ}\ \text{」}.
\]
注意若 $X_*(T)_\Gamma$ 有挠，则该映射不是单射。
\end{proposition}

\begin{proof}[略证]
\srcpage{57}先处理 $T=\mathbb{G}_m$ 的情形。此时 $B(T)=\breveE^\times/(\sigma-1)$，有满射
$\twoheadrightarrow\ZZ$，$b\mapsto\nu(b)$。容易看出这是同构；事实上它就是秩 $1$
等晶体分类的重新表述。

再处理 $T=\Res_{E'/E}\mathbb{G}_m$ 的情形，其中 $E'/E$ 是有限可分扩张。则由 Shapiro
引理的变体有
\[
B(E,T)\cong B(E',\mathbb{G}_m)\cong\ZZ,\qquad
B(E,\Res_{E'/E}G)\cong B(E',G).
\]
同样由 Shapiro 引理，$X_*(T)\cong\Ind_{\Gamma_{E'}}^{\Gamma_E}(\ZZ)$，且
$X_*(T)_{\Gamma_E}\cong\ZZ_{\Gamma_{E'}}\cong\ZZ$。可以直接验证这些等同与上述取平均的
映射相容（事实上，那个映射本来就可以定义为使它们相容的映射）。

最后，用诱导环面来分解一般的环面。任何环面 $T$ 都承认满射
\[
\prod_{i=1}^n\Res_{E_i/E}\mathbb{G}_m\twoheadrightarrow T.
\]
这化到前一种情形。但要指出，正是在这一步会出现挠。
\end{proof}

回到一般的 $G$。可以定义「Borovoi 基本群」
\[
\Gamma\curvearrowright\pi_1(G)=\pi_1(G_{\bar E})=X_*(T)/\text{余根格}.
\]
要点在于：对 $G/\CC$，这会还原出通常的拓扑 $\pi_1$（而不是投射有限的东西！）。

\begin{proposition}
存在唯一的函子性延拓
\[
\kappa\colon B(G)\to\pi_1(G)_\Gamma,
\]
它延拓环面上的映射 $B(T)\xrightarrow{\sim}X_*(T)_\Gamma=\pi_1(T)_\Gamma$。
\end{proposition}

\begin{proof}[略证]
先考虑 $G_{\mathrm{der}}$ 单连通的情形。此时有短正合列
\[
1\to G_{\mathrm{der}}\to G\to\underbrace{D}_{\text{环面}}\to1,
\]
且 $\pi_1(G)\xrightarrow{\sim}\pi_1(D)$，故 $\kappa$ 由投影到 $D$ 定义。

对一般的 $G$，存在 z-扩张 $G'\twoheadrightarrow G$（中心核），使
$G'_{\mathrm{der}}$ 单连通。于是 $B(G')\twoheadrightarrow B(G)$，我们试图定义
\begin{center}
\begin{tikzcd}
B(G') \arrow[r, two heads] \arrow[d] & B(G) \arrow[d, dashed, "\exists"] \\
\pi_1(G')_\Gamma \arrow[r] & \pi_1(G)_\Gamma
\end{tikzcd}
\end{center}
\end{proof}

\begin{example}
取 $E=\Fq((t))$。此时把 $G(\bar{\Fq}((t)))$ 想成「$G$ 中的代数环路」。存在映射
$G(\bar{\Fq}((t)))\to\pi_1(G)_\Gamma$，它可以理解为记住该类的「拓扑环路」。它下降为
Kottwitz 映射。
\begin{center}
\begin{tikzcd}
G(\bar{\Fq}((t))) \arrow[d] \arrow[dr] & \\
B(G) \arrow[r, "\kappa"] & \pi_1(G)_\Gamma
\end{tikzcd}
\end{center}
\end{example}

\begin{theorem}[Kottwitz]
\srcpage{58}对所有的 $G$，映射
\[
(\nu,\kappa)\colon B(G)\to(X^+_{\QQ})^\Gamma\times\pi_1(G)_\Gamma
\]
是单射。
\end{theorem}

由此可在 $B(G)$ 上定义偏序：若 $\nu(b)\prec\nu(b')$（支配序）且 $\kappa(b)=\kappa(b')$，
则称 $b\le b'$。该序下的极小元称为基本的。

\begin{proposition}
Kottwitz 映射 $\kappa$ 诱导 $B(G)^{\mathrm{basic}}\xrightarrow{\sim}\pi_1(G)_\Gamma$。
\end{proposition}

非基本的元素可以用 Levi 子群来理解（至少在 $G$ 拟分裂时）。

\begin{proposition}
元素 $b\in B(G)$ 是基本的，当且仅当 $\nu(b)$ 是中心的。
\end{proposition}

对任何 $b\in B(G)$，可以看 $b$ 的 $\sigma$-中心化子，它等同于相应 $\otimes$-函子的
自同构。这定义了 $E$ 上一个连通约化群 $G_b$。若 $b$ 是基本的，则 $G_b$ 是 $G$ 的内形式。
更一般地，若 $G$ 拟分裂，则它是 $G$ 的某个 Levi 子群（即 $\nu(b)$ 的中心化子）的内形式。
这个群通常记作 $J_b$，不过对 $b=1$ 我们宁愿写 $G$ 而不写 $J_1$。

\subsection{Fargues--Fontaine 曲线上的 $G$-丛}

固定通常的记号：$E\supset\cO_E\ni\pi$，剩余域为 $\Fq$。取定 $\bar{\Fq}$，于是
$\breveE=W_{\cO_E}(\bar{\Fq})[1/\pi]$。

\begin{definition}
设 $S\in\Perf_{\Fq}$。$X_S$ 上的 $G$-主丛是正合 $\otimes$-函子
\[
\cE\colon\Rep_E(G)\to\mathrm{VB}(X_S).
\]
\end{definition}

\begin{definition}
称 $\Bun_G$ 为 Fargues--Fontaine 曲线上 $G$-丛的叠，它是 $\Perf_{\Fq}$ 上的 v-叠，把
\[
S\mapsto\{\text{$X_S$ 上的 $G$-丛}\}.
\]
（它是 v-叠这一事实来自向量丛的 v-下降。）
\end{definition}

\begin{warning}
不存在「那个」Fargues--Fontaine 曲线，因为 $X_{C,E}$ 的构造依赖于某个输入域 $C$。
但「Fargues--Fontaine 曲线上 $G$-丛的叠」是良定义的。
\end{warning}

\begin{theorem}[{Fargues（当 $E/\Qp$ 时），Anschütz（一般情形）}]
若 $S=\Spa(C,C^+)$，其中 $C$ 是完备代数闭域，则把 $\Spa\breveE/\sigma^\ZZ$ 上的
$G$-主丛送到它在 $Y_S/\varphi^\ZZ=X_S$ 上的拉回的函子
\[
G\text{-}\Isoc\to\Bun_G(S)
\]
在同构类上诱导双射
\[
\Bun_G(S)/\sim\ \xrightarrow{\sim}\ B(G),
\]
它甚至是同胚 $|\Bun_G(S)|\cong B(G)$。
\end{theorem}

\begin{proof}[略证]
设 $\cE$ 是 $X_C:=X_S$ 上的 $G$-主丛。对任何 $V\in\Rep_E(G)$，$\cE(V)$ 带有
HN-滤过，于是得到函子
$\Rep_E(G)\to\QQ\text{-}\Fil\mathrm{VB}(X_C)^{\mathrm{HN}}$，即 $X_C$ 上 $\QQ$-滤过的
向量丛的范畴，其中所有 $\cE^\lambda$ 都是斜率为 $\lambda$ 的半稳定丛。

这仍是 $\otimes$-函子，可由向量丛的分类验证。（「张量积上的 HN 滤过是各 HN 滤过之
张量积」这一断言一般很难，但在本题中由分类定理很容易。）还需要得到正合性。当
$E/\Qp$ 时这是平凡的，因为此时 $\Rep_E(G)$ 半单，问题只是加性。在正特征情形则更微妙：
因为 $\Rep_E(G)$\srcpage{59} 很复杂；我们改用 Haboush 证明的某些「几何约化性」作为完全可约性的替代。

现在取分次化，得到映到 $\QQ\text{-}\mathrm{GrVB}(X_C)^{\mathrm{HN}}$ 的投影。它实际上
就等价于 $\Isoc_E$，因为半稳定向量丛的子范畴恰与等晶体的范畴等价。这给出一个候选的
$G$-等晶体，剩下要分裂滤过。为此用到 $H^1(X_C,\cO(\lambda))=0$（$\lambda>0$）。
\end{proof}

\begin{corollary}
等晶体 $b\in B(G)$ 是基本的，当且仅当 $\cE_b\in\Bun_G(X_C)$ 在 Atiyah--Bott 意义下
半稳定。
\end{corollary}

\begin{theorem}
映射 $|\Bun_G|\to B(G)$ 连续，即
\[
-\nu\colon|\Bun_G|\to(X^+_{\QQ})^\Gamma\ \text{是上半连续的},
\]
并且
\[
-\kappa\colon|\Bun_G|\to\pi_1(G)_\Gamma\ \text{是局部常值的}.
\]
事实上，
\[
\kappa\colon\pi_0(\Bun_G)\xrightarrow{\sim}\pi_1(G)_\Gamma.
\]
证明将在下一次给出。
\end{theorem}

\begin{remark}
这是 Rapoport--Richartz [RR96] 关于 $G$-等晶体族的定理的类比（不过特化箭头的方向相反）。
他们没能在完全一般的情形下证明 $\kappa$ 局部常值这一断言。本定理可用来补全他们的结果。
要点在于它给出了 Kottwitz 的映射 $\kappa$ 的一个漂亮的几何描述，而该映射先验地由一个
相当复杂的过程定义。
\end{remark}
